% ==================================================
% 5. PRODUCT FEATURES
% ==================================================


\section{Product Features}

\PAChange{This section describes the main
functional user requirements of the
Medication Safety Assistant as high-level
product features. Each feature explains
what the system provides and the value
it offers to its users.}

\PAChange{Detailed user--system interactions
are specified in the Use-case Specifications.
Planned extensions are separated from
the current prototype scope.}



% ==================================================
% PF-01: MULTIMODAL MEDICATION INTAKE
% ==================================================

\subsection{PF-01: Multimodal Medication Intake}

Users can provide medication information
through four input methods: prescription-image
OCR, NLP Chatbot, Voice/STT, and Manual Input.

The system supports information extraction,
medication search suggestions, and user
corrections. Users can switch to another
input method when needed.

\PAChange{The system records available medication
details and presents them for user review.
Failed or incomplete input processing
allows correction, retry, or manual entry.}


% ==================================================
% PF-02: INTENT RECOGNITION AND USER CONFIRMATION
% ==================================================

\subsection{PF-02: Intent Recognition and User Confirmation}

\PAChange{During medication intake, the system
identifies the requested task through UI
selection or NLP analysis of text and voice
input. Unclear requests require clarification.}

\PAChange{Users must review, correct, and confirm
recorded medication information before
safety checks or reminder activation.
Confirmation verifies the recorded input,
not medication identity or medical safety.}

% ==================================================
% PF-03: VIETNAMESE MEDICATION IDENTIFICATION
% ==================================================

\subsection{PF-03: Vietnamese Medication Identification}

The system identifies Vietnamese medicines
and their active ingredients using available
medication data sources.

It supports matching different medicine
names to their active ingredients and
related medication information.
Unknown or unclear matches are reported
instead of being treated as confirmed.

\PAChange{Ambiguous medicines may require
user clarification. Unresolved medicines
are not used in checks that require
confirmed ingredient information.}


% ==================================================
% PF-04: MEDICATION RISK ASSESSMENT
% ==================================================

\subsection{PF-04: Medication Risk Assessment}

\PAChange{The system checks supported
drug--drug interactions using available
DDInter 2.0 data and verified ingredient
mappings. It also checks for duplicate
active ingredients when sufficient
information is available.}

It identifies potential risks, their
severity when available, and cases
that need further review. Missing data
or no detected interaction does not
guarantee that a medication plan is safe.

\PAChange{Assessment results are marked
as assessed, partially assessed,
unassessed, or failed. Severe interactions
trigger immediate warnings and the
new-reminder safety gate.}

The assessment can be extended to other
medication risks when reliable data
and suitable methods become available.


% ==================================================
% PF-05: SAFETY RESULTS AND RECOMMENDATIONS
% ==================================================

\subsection{PF-05: Safety Results and Recommendations}

The system presents medication safety
results in a clear and understandable way,
including detected risks, involved medicines,
available explanations, and data limitations.

It provides appropriate next steps and
recommends professional review when needed.
The system supports informed decisions
without independently changing prescribed
treatment.

\PAChange{Explanations use supported data.
Results distinguish detected risks from
unassessed checks. Finding no issue does
not prove that all medicines are safe.}


% ==================================================
% PF-06: HEALTHCARE PROFESSIONAL REVIEW
% ==================================================

\subsection{PF-06: Healthcare Professional Review}

\PAChange{Patients can request HCP review
of a saved safety case and permit access
to selected case information. Authorized
HCPs can review safety findings, patient
summaries, and data limitations.}

They can request additional information,
provide feedback, and record review decisions.
Cases awaiting required review remain pending.

\PAChange{Patients and HCPs can exchange
messages and follow-up questions within
the same review case. The system keeps
the conversation, review outcomes, and
request status for later reference.}

\PAChange{Cases waiting for patient information
or HCP responses remain clearly marked.
Only recorded, eligible HCP confirmations
can affect pending new reminders.
All HCP reviews are simulated and do not
represent real clinical approval.}


% ==================================================
% PF-07: MEDICATION REMINDER MANAGEMENT
% ==================================================

\subsection{PF-07: Medication Reminder Management}

\PAChange{Users can create reminders from
existing prescriptions or recorded HCP
instructions after confirming the schedule.
Unverified schedules cannot be activated.}

\PAChange{Users can change reminder
notification settings without changing
prescribed treatment. Proposed changes
to medication timing require simulated
HCP confirmation before activation.}

\PAChange{If a severe drug interaction
is detected, new reminders for affected
medicines remain pending HCP review.
Unrelated eligible reminders may proceed.
Existing reminders and prescribed medicines
are never stopped automatically.}


% ==================================================
% PF-08: REMINDERS AND ADHERENCE TRACKING
% ==================================================

\subsection{PF-08: Medication Reminders and Adherence Tracking}

The system provides medication reminders
and allows users to record doses as
taken, missed, or skipped.

Users can review reminder schedules
and medication adherence records.
This information helps users track
their medication adherence and may
support future adherence analysis.


\PAChange{Reminder notifications depend on
browser support and user permissions.
Adherence records include relevant dates
and times and do not automatically
change prescribed treatment.}


% ==================================================
% PF-09: MEDICATION AND HEALTH HISTORY
% ==================================================


\subsection{PF-09: Medication and Health History}

\PAChange{The system stores confirmed medication
records, previous safety assessments,
detected risks, and HCP review outcomes,
including their dates and statuses.}

\PAChange{Patients can view their saved
medication information and history.
New prescriptions are recorded through
the confirmed medication intake process.
Corrections must match the provided
source and require user confirmation.}

\PAChange{Previous records and assessment results
remain available as history. New information
does not automatically replace prescribed
treatment or existing reminders.
HCP access to patient history requires
authorization.}




% ==================================================
% PF-10: CONTEXT-AWARE CHATBOT AND PATIENT SUPPORT
% ==================================================

\subsection{PF-10: Context-Aware Chatbot and Patient Support}

\PAChange{The AI chatbot supports multi-turn
medication-related requests using supported
reference data and permitted patient information.
It is separate from the patient--HCP review
conversation described in PF-06.}

\PAChange{It asks follow-up questions
when needed and reports limitations
when reliable information is unavailable.
It does not independently diagnose,
prescribe, or change treatment.}

\PAChange{The system can generate structured
patient health summaries from available
records. Patients review the summaries
before saving or sharing them with
authorized HCPs. Sharing requires
patient permission.}

\PAChange{\textbf{Planned extensions:}
Consultation preparation, medication
record comparison, and explanations
of confirmed medication instructions.}


% ==================================================
% PF-11: MEDICATION DATASET REVIEW
% ==================================================

\subsection{PF-11: Medication Dataset Review}

\PAChange{Dataset Reviewers can inspect
medication records, active ingredients,
mappings, and supporting sources
through a restricted interface.}

\PAChange{They can record data issues and
review decisions such as Accept, Reject,
or Needs Correction. The system keeps
source references and review history.}

\PAChange{Unverified records must not
be presented as clinically validated.
Dataset review does not replace
professional clinical validation.}


% ==================================================
% PF-12: USER AND ACCESS MANAGEMENT
% ==================================================

\subsection{PF-12: User and Access Management}

\PAChange{Registered users can sign in
to access features permitted by their
assigned roles. Administrators can manage
user accounts, account status, roles,
and access permissions.}

\PAChange{Role-Based Access Control (RBAC)
limits access to authorized features
and information for the five defined
user roles. Relevant administrative
activities are recorded for review.}


% ==================================================
% PF-13: TECHNICAL SUPPORT
% ==================================================

\subsection{PF-13: Technical Support}

\PAChange{Users can report technical
issues. Technical Support personnel
can review reported problems, investigate
permitted diagnostic information,
record solutions, and track issue status.}

\PAChange{Support access is limited
to information needed for troubleshooting.
Private patient records are protected
from unnecessary access.}



% ==================================================
% PRODUCT FEATURE AND USE-CASE TRACEABILITY
% ==================================================

\subsection{Product Feature and Use-Case Traceability}

\PAChange{The following table links the
13 product features to the 18 use cases
in the current Use-case Specifications.
A use case may support multiple features.
Detailed flows are defined in the
Use-case Specifications.}

\begin{longtable}{|
    p{0.18\textwidth}|
    p{0.72\textwidth}|
}
\hline

\textbf{Feature} &
\textbf{Related Use Cases}
\\ \hline

\endfirsthead

\hline
\textbf{Feature} &
\textbf{Related Use Cases}
\\ \hline

\endhead

\PAChange{PF-01}
& \PAChange{UC-02: Enter and Confirm
Medication Information}
\\ \hline

\PAChange{PF-02}
& \PAChange{UC-02: Enter and Confirm
Medication Information}
\\ \hline

\PAChange{PF-03}
& \PAChange{UC-03: Check Medication Safety}
\\ \hline

\PAChange{PF-04}
& \PAChange{UC-03: Check Medication Safety}
\\ \hline

\PAChange{PF-05}
& \PAChange{UC-03: Check Medication Safety;
UC-04: Request HCP Review}
\\ \hline

\PAChange{PF-06}
& \PAChange{UC-04: Request HCP Review;
UC-05: Review Medication Safety Case;
UC-13: Share Patient Health Summary with HCP}
\\ \hline

\PAChange{PF-07}
& \PAChange{UC-03: Check Medication Safety;
UC-05: Review Medication Safety Case;
UC-06: Set Up Medication Reminders;
UC-07: Manage Medication Reminders}
\\ \hline

\PAChange{PF-08}
& \PAChange{UC-06: Set Up Medication Reminders;
UC-07: Manage Medication Reminders;
UC-08: Record Dose Status;
UC-09: View Adherence History}
\\ \hline

\PAChange{PF-09}
& \PAChange{UC-02: Enter and Confirm
Medication Information;
UC-03: Check Medication Safety;
UC-05: Review Medication Safety Case;
UC-10: View Medication and Safety History}
\\ \hline

\PAChange{PF-10}
& \PAChange{UC-11: Ask a Medication Question;
UC-12: Generate Patient Health Summary;
UC-13: Share Patient Health Summary with HCP}
\\ \hline

\PAChange{PF-11}
& \PAChange{UC-14: Review Medication Dataset Record}
\\ \hline

\PAChange{PF-12}
& \PAChange{UC-01: Sign In;
UC-15: Manage User Accounts and Roles;
UC-16: View Administrative Activity History}
\\ \hline

\PAChange{PF-13}
& \PAChange{UC-17: Report Technical Issue;
UC-18: Handle Technical Support Request}
\\ \hline

\end{longtable}
